% ============================================================
% Preámbulo común a todos los documentos del Proyecto 2
% ============================================================
\usepackage[utf8]{inputenc}
\usepackage[T1]{fontenc}
\usepackage[spanish]{babel}
\usepackage[margin=2.3cm]{geometry}
\usepackage{amsmath}
\usepackage{booktabs}
\usepackage{float}
\usepackage{array}
\usepackage{enumitem}
\usepackage{tikz}
\usetikzlibrary{shapes.geometric, arrows.meta, positioning}
\usepackage[hidelinks]{hyperref}

\renewcommand{\contentsname}{Índice}
\renewcommand{\figurename}{Figura}
\renewcommand{\tablename}{Tabla}

\setlist{itemsep=2pt, topsep=4pt}

\tikzset{
  bloque/.style={rectangle, rounded corners, draw=black, thick, fill=blue!8,
                 text width=7.5cm, align=center, font=\small, inner sep=6pt},
  bloqueMCU/.style={bloque, fill=green!12},
  bloqueSPI/.style={bloque, fill=orange!25},
  bloqueSBC/.style={bloque, fill=violet!12},
  flujo/.style={-{Stealth[length=2.5mm]}, thick},
  bidir/.style={{Stealth[length=2.5mm]}-{Stealth[length=2.5mm]}, thick},
  inicio/.style={rectangle, rounded corners=6pt, draw=black, thick, fill=green!15,
                 minimum width=3.5cm, minimum height=0.8cm, align=center, font=\small},
  proc/.style={rectangle, draw=black, thick, fill=blue!8, text width=5.2cm,
               align=center, font=\small, minimum height=0.9cm, inner sep=4pt},
  dec/.style={diamond, aspect=2.4, draw=black, thick, fill=yellow!20,
              text width=3.0cm, align=center, font=\footnotesize, inner sep=1pt},
}

% Datos del proyecto (un solo lugar para mantenerlos)
\newcommand{\Universidad}{Universidad Nacional de Colombia --- Sede Manizales}
\newcommand{\Facultad}{Facultad de Ingeniería y Arquitectura}
\newcommand{\Carrera}{Carrera de Ingeniería Electrónica}
\newcommand{\Materia}{Programación de sistemas Linux embebido}
\newcommand{\TituloProyecto}{Sistema Híbrido de Adquisición Analógica mediante SPI\\[0.3cm]y Dashboard Embebido en Linux}
\newcommand{\FechaDoc}{Manizales, noviembre de 2026}

% Comando de color-neutral para valores a completar por el equipo
\newcommand{\completar}[1]{\textit{[#1]}}